%% Rechnungsvorlage – wird vom InvoiceRenderer gefüllt.
%% Platzhalter haben die Form <<...>> und werden vollständig ersetzt. Welche es
%% gibt, zählt TemplateValue auf. Die Vorlage muss nicht jeden setzen, aber
%% jeden, der an Rechnungsart oder Rechnung hängt (#188).
%%
%% Jede einzelne Angabe kommt als eigener Platzhalter und darf leer sein. Ob sie
%% dasteht, entscheidet der Renderer, indem er sie leer lässt oder nicht; wo und
%% mit welcher Beschriftung, entscheidet diese Vorlage. Sie fragt dafür nur
%% \ifblank und nichts sonst: Was an der Rechnungsart hängt, soll nicht in
%% jedem Design neu entschieden werden, ohne Compiler.
%%
%% Gestaltungsprinzip: EIN Blickfang (der Titelbalken), sonst ruhiger Textfluss.
%% Linien gibt es nur dort, wo sie Zahlen gliedern – nicht als Seitenraster.
%% Gesperrte Versalien (das Stilmittel des Lebenslaufs) kommen genau einmal vor,
%% als Spaltenkopf der Positionstabelle.
\documentclass[11pt,a4paper]{article}

%% pdfx erzeugt das PDF/A-1b: XMP-Metadaten, OutputIntent mit sRGB-Profil und
%% eingebettete Schriften. Muss früh geladen werden, es zieht hyperref mit.
%% Die Metadaten kommen aus der gleichnamigen .xmpdata-Datei daneben.
\usepackage[a-1b]{pdfx}

\usepackage{fontspec}
\usepackage{geometry}
%% xcolor lädt pdfx bereits selbst – ein zweites \usepackage mit Optionen
%% gäbe einen "Option clash". colortbl direkt nachladen liefert \arrayrulecolor.
\usepackage{colortbl}
\usepackage{booktabs}  %% \midrule bringt die vertikalen Abstände selbst mit
\usepackage{tabularx}
\usepackage{array}
\usepackage{etoolbox}  %% \ifblank für die Angaben, die leer sein dürfen

%% Dieselbe Schrift wie im Lebenslauf – dadurch passen die Dokumente zusammen.
%%
%% Ligatures=TeXOff: fontspec zieht für die Hauptschrift sonst Folgen wie
%% "--", "<<" oder zwei Hochkommas zu einem Zeichen zusammen (Strich, Guillemet,
%% Anführungszeichen). Die Vorlage braucht das nicht, sie schreibt keine dieser
%% Folgen; getroffen hätte es nur Text aus den Eingabedaten, und der soll so
%% dastehen, wie er in der YAML-Datei steht. In escapeLatex aufbrechen geht
%% nicht: Nachgemessen trennt LuaLaTeX die Ligatur weder bei "-{}-" noch bei
%% "{<}{<}". Die Schriftligaturen (fi, fl) bleiben, sie ändern keinen Text.
\setmainfont{texgyreheros}[
  Ligatures      = TeXOff,
  Extension      = .otf,
  UprightFont    = *-regular,
  BoldFont       = *-bold,
  ItalicFont     = *-italic,
  BoldItalicFont = *-bolditalic,
]

%% Farben von fluch-it-consulting.de
%% Mittelweg zwischen dem Rot der Webseite (#8B1A2B) und dem gedämpften Ton
%% der bisherigen Rechnung (#8E4848): kanalweise das arithmetische Mittel.
%% Beide Ausgangsfarben stehen daneben, das Umschalten ist eine Zeile.
\definecolor{akzent}{HTML}{8D313A}
\definecolor{akzentweb}{HTML}{8B1A2B}
\definecolor{akzentalt}{HTML}{8E4848}
\definecolor{haupttext}{HTML}{1C1C1E}
\definecolor{gedaempft}{HTML}{5C5C62}
\definecolor{linie}{HTML}{C9C9CF}

%% Tabellenlinien global. Muss AUSSERHALB der Tabelle stehen: Mitten im
%% Tabellenkörper eröffnet \arrayrulecolor eine neue Zelle, und das \noalign
%% von \midrule landet dann an der falschen Stelle.
\arrayrulecolor{linie}

\geometry{a4paper, left=25mm, right=20mm, top=18mm, bottom=16mm}

\color{haupttext}
\setlength{\parindent}{0pt}
\pagestyle{empty}

%% Flattersatz statt Blocksatz. In den schmalen Spalten einer Rechnung reißt
%% Blocksatz sichtbare Löcher zwischen die Wörter; Heros verträgt Flattersatz gut.
\raggedright

%% Keine Silbentrennung: Auf einer Rechnung stehen fast nur Eigennamen,
%% Beträge und Adressen – getrennte Firmennamen sehen schlicht falsch aus.
\hyphenpenalty=10000
\exhyphenpenalty=10000
\tolerance=2000

%% Satzabstand: Nach Punkt, Doppelpunkt und Komma steht der gewöhnliche
%% Wortzwischenraum, nicht der breitere englische (#143). Von sich aus hielte
%% TeX jeden Punkt nach einem Kleinbuchstaben oder einer Ziffer für ein
%% Satzende – "Dr. Erika", "14. August", "USt-IdNr. DE" –, und aus den
%% Eingabedaten lässt sich nicht ablesen, wo ein Satz wirklich endet. Den
%% breiten Abstand bekommen nur die eigenen Sätze des Programms, siehe
%% \hinweisabsatz.
\frenchspacing

%% Achtung: minipage ruft \@parboxrestore auf und setzt damit die
%% Absatzparameter zurück – der Flattersatz von oben gilt dort NICHT.
%% Ein \AtBeginEnvironment-Hook hilft nicht, weil \@parboxrestore erst danach
%% läuft. Deshalb steht \raggedright unten in jeder Minipage mit Fließtext.

%% Kleine, gesperrte Versalien – dasselbe Stilmittel wie im Lebenslauf.
%% \rubrik statt \label, das ist in LaTeX vergeben.
\newlength{\balkenhoehe}

%% Abstände als benannte Längen (#51): Vorher standen die Millimeter in Kotlin –
%% wer das Satzbild ändern wollte, musste zwei Dateien anfassen und beide treffen.
%%
%% 2.5mm zwischen Positionen statt 1.5mm: Sobald eine Position eine Zusatzzeile
%% hat, verschmilzt sie bei knappem Abstand optisch mit der nächsten. Dasselbe
%% gilt für die abgesetzten Abschläge, die immer eine Zusatzzeile haben.
\newlength{\positionsabstand}\setlength{\positionsabstand}{2.5mm}
\newlength{\summenabstand}\setlength{\summenabstand}{1mm}
\newlength{\blockabstand}\setlength{\blockabstand}{3mm}
\newlength{\ustidabstand}\setlength{\ustidabstand}{2mm}

%% Die Zeilenwechsel, die diese Abstände tragen (#76). Der Renderer ruft nur das
%% Makro auf und kennt die Länge nicht. Vorher schrieb er \\[\positionsabstand]
%% selbst, und ein vertippter Längenname fiel erst LuaLaTeX auf – ein Makro prüft
%% das Setzen, eine Länge nicht. \blocktrenner und \ustidtrenner ruft seit #188
%% nur noch diese Vorlage selbst auf.
\newcommand{\positionstrenner}{\\[\positionsabstand]}
\newcommand{\summentrenner}{\\[\summenabstand]}
\newcommand{\blocktrenner}{\\[\blockabstand]}
\newcommand{\ustidtrenner}{\\[\ustidabstand]}

%% Der schlichte Zeilenwechsel ohne Abstand: zwischen den Zeilen der Anschrift,
%% der Steuerhinweise und im Briefkopf (#147).
%% Eines für alle Blöcke. raster kommt damit ebenso aus (ADR-0011); braucht ein
%% Design für verschiedene Blöcke verschiedene Zeilenwechsel, kostet das einen
%% Eintrag in TemplateMacro.
%%
%% Aufgerufen wird es als \zeilenwechsel{}, vom Renderer wie von dieser Vorlage,
%% und das leere Paar schützt die nächste Zeile: Beginnt sie mit "[", läse \\
%% sie sonst als Abstandsangabe. Ein rohes \\ steht deshalb nur vor festem
%% Text, nie vor einem Platzhalter – auch nicht über ein Zeilenende hinweg,
%% denn \\ überliest den Leerraum bis zum nächsten Zeichen.
\newcommand{\zeilenwechsel}{\\}

%% Die Person unter der Firmierung im Briefkopf (#147). Das {} hinter \\ aus
%% demselben Grund wie beim Aufruf von \zeilenwechsel.
\newcommand{\kopfperson}[1]{\\{}#1}

%% Eine Zeile im Block oben rechts: Beschriftung und Wert. Ist der Wert leer,
%% entfällt die ganze Zeile, samt Beschriftung – sonst stünde auf einer
%% Stornorechnung "Zahlbar ohne Abzug bis" ohne Datum.
%%
%% Jede Zeile beendet sich selbst mit \par. Das ist, was \\ in einem Block mit
%% \raggedleft ohnehin tut (latex.ltx, \@centercr); ein \\ hinter der letzten
%% Zeile ließe dagegen eine leere Zeile zurück.
\newcommand{\kopfangabe}[2]{\ifblank{#2}{}{#1 #2\par}}

%% Diese Angabe bricht nicht um. Im schmalen Block oben rechts bräche
%% "Juni bis Oktober 2026" sonst hinter dem "bis" und läse sich wie zwei Angaben.
\newcommand{\zusammen}[1]{\mbox{#1}}

\newcommand{\rubrik}[1]{{\footnotesize\color{gedaempft}\addfontfeature{LetterSpace=14}\MakeUppercase{#1}}}

%% Ein Hinweisabsatz unter den Summen – der Pflichtsatz der Rechnungsart oder
%% der Satz zum ausgeglichenen Restbetrag. Er bringt seinen Abstand selbst mit,
%% weil er nur bei einzelnen Rechnungsarten gesetzt wird: Stünde der \vspace in
%% der Vorlage, nähme er bei allen anderen Raum ein, obwohl dort nichts steht.
%%
%% Er trägt die eigenen Sätze des Programms, und nur in ihm steht nach einem
%% Satzende der breitere Abstand (#143): nach . ? !, aber nicht nach : ; , –
%% das sind nie Satzenden. Die Gruppe hält das auf das Argument beschränkt.
%% Eingabedaten gehören deshalb nicht hier hinein: Eine Abkürzung wie "Dr."
%% bekäme sonst wieder den breiten Abstand. Das % hinter \noindent verhindert
%% einen Leerraum am Absatzanfang.
\newcommand{\hinweisabsatz}[1]{\par\vspace{6mm}\noindent%
  {\nonfrenchspacing\sfcode`\:=1000 \sfcode`\;=1000 \sfcode`\,=1000 #1}}

%% Der Kopf des Absetzungsblocks: wie ein Hinweisabsatz, aber mit Luft zur
%% Tabelle darunter – sonst klebt die erste Rechnung an der Ankündigung.
\newcommand{\absetzungskopf}[1]{\par\vspace{6mm}\noindent #1\par\vspace{1.5mm}}

%% Die Bankverbindung mitsamt ihrer Überschrift. Als Makro und nicht als fester
%% Text, weil sie nur auf Dokumenten steht, die zur Zahlung auffordern: Eine
%% Stornorechnung kündigt eine Erstattung an, und eine Überschrift ohne Konto
%% darunter wäre schlimmer als keine. Ob es dasteht, entscheidet der Renderer,
%% indem er es aufruft oder nicht; was darin steht, diese Vorlage. Der Abstand
%% darüber gehört deshalb ins Makro, wie bei \hinweisabsatz.
\newcommand{\bankblock}{\vspace{9mm}\noindent{\bfseries Bankverbindung}\\[1.5mm]%
  Kontoinhaber: <<KONTOINHABER>>\\
  Institut: <<INSTITUT>>\\
  IBAN: <<IBAN>>\\
  BIC: <<BIC>>}

%% Die Tabelle der abgesetzten Abschläge: links die Rechnung, rechts der Betrag.
%% Dieselbe Spaltenaufteilung wie beim Summenblock, nur über die volle Breite –
%% die Zeilen tragen Nummer und zwei Daten und passen nicht in 0,55.
\newcommand{\absetzungstabelle}[1]{%
  \begin{tabularx}{\linewidth}{@{}>{\raggedright\arraybackslash}X r@{}}#1\end{tabularx}}

%% Die Aufteilung eines Abschlags in Netto und Steuer: klein, gedämpft und
%% eingerückt unter der Zeile, zu der sie gehört.
\newcommand{\absetzungsdetail}[1]{\newline\hspace*{4mm}{\footnotesize\color{gedaempft}#1}}

%% Dasselbe für eine Position, nur ohne Einzug – die Beschreibungsspalte beginnt
%% ohnehin am linken Rand. Zwei Makros und nicht eines: Der Unterschied im Einzug
%% ist Satzbild und wird nicht nebenbei in einem Umbau eingeebnet (#51).
\newcommand{\positionsdetail}[1]{\newline{\footnotesize\color{gedaempft} #1}}

%% Eine Zeile der Positionstabelle. Damit steht die Zahl der Spalten nur noch
%% hier: Vorher zählte die Vorlage fünf und der Renderer verband fünf Zellen mit
%% kaufmännischem Und, ohne dass etwas die beiden Zahlen verglich.
%%
%% Das erste Argument ist die laufende Nummer der Position. Dieses Design setzt
%% sie nicht.
\newcommand{\positionszeile}[6]{#2 & #3 & #4 & #5 & #6}

%% Eine Zeile im Summen- oder Absetzungsblock: Bezeichnung links, Betrag rechts.
\newcommand{\summenzeile}[2]{#1 & #2}

%% Die fette Zeile unter dem Strich, mit der ein Betragsblock endet – der
%% Gesamtbetrag der Summen, der Restbetrag der Absetzung.
%%
%% Sie beendet selbst die Zeile davor und steht deshalb nur hinter einer Zeile,
%% nie als erste (#147). Ein Trenner vor ihr ginge nicht: Der Renderer schreibt
%% ihn als \name{}, und ein {} vor \midrule ist ein "Misplaced \noalign".
\newcommand{\abschlusszeile}[2]{\\\midrule\textbf{#1} & \textbf{#2}}

\begin{document}

%% ------------------------------------------------------------------ Kopf
%% Alles rechtsbündig in einem Block, wie auf der bisherigen Rechnung: die
%% Firmierung fett, darunter Anschrift und Kontakt in derselben Schriftgröße.
%% Links bleibt frei – der Balken weiter unten ist der einzige Blickfang.
%%
%% Die Person steht unter der Firmierung nur, wenn es keine Geschäftsführung
%% gibt: Die steht rechts neben dem Empfänger, und bei einer UG sind beide
%% meist dieselbe Person. Für Einzelunternehmen bleibt der Name hier stehen,
%% §14 Abs. 4 Nr. 1 UStG verlangt ihn. Die Webseite steht bewusst nicht im
%% Kopf – sie wiederholte nur die Domain der E-Mail. Die E-Mail gehört hierher:
%% Rückfragen zu einer Rechnung kommen meist schriftlich.
{\raggedleft\small\color{gedaempft}
  {\bfseries\color{haupttext}<<FIRMA>>}\ifblank{<<GESCHAEFTSFUEHRUNG>>}{\kopfperson{<<INHABER>>}}{}\zeilenwechsel{}%
  <<STRASSE>> · <<PLZ>> <<ORT>>\ifblank{<<TELEFON>>}{}{\zeilenwechsel{}Tel. <<TELEFON>>}\zeilenwechsel{}<<EMAIL>>\par}

\vspace{18mm}

%% ----------------------------------------------------------- Empfängerfeld
%% Ohne Rücksendezeile: Die Rechnung wird als PDF versandt, nie im
%% Fensterumschlag. Rechts daneben die Vertretung und die Angaben zur
%% Rechnung. 0.44 der Breite = 72,6 mm; nachgemessen braucht die längste
%% Zeile ("Leistungszeitraum: 01.07.2026 – 31.08.2026") 69,8 mm.
%%
%% Die USt-IdNr. des Kunden unter seiner Anschrift (BT-48): Bei Reverse Charge
%% und innergemeinschaftlicher Lieferung verlangt §14a UStG die USt-IdNr. beider
%% Seiten, die eigene steht im Fuß. Sie hängt ohne Leerraum an der Anschrift,
%% sonst stünde ohne sie ein Leerzeichen am Ende des Blocks.
\begin{minipage}[t]{0.52\linewidth}\raggedright
  <<EMPFAENGER>>\ifblank{<<KUNDE_USTID>>}{}{\ustidtrenner{}USt-IdNr. <<KUNDE_USTID>>}
\end{minipage}%
\hfill
%% Die Geschäftsführung steht hier und nicht im Fuß, sonst bräche die erste
%% Fußzeile um. "Geschäftsführung" statt "Geschäftsführer": Die Bezeichnung
%% steht neben einem Namen, dessen Geschlecht die Stammdaten nicht kennen.
%%
%% Nummer und Datum der Rechnung stehen im Balken, nicht hier. Die
%% Kundennummer steht als erste Kopfangabe, gleich neben dem Kunden. Die Beschriftung
%% des Zeitraums kommt vom Renderer, weil sie an der Rechnungsart hängt
%% ("Voraussichtlicher Leistungszeitraum", UStAE 14.8). Der Zeitraum bricht
%% nicht um: Im schmalen Block bräche "Juni bis Oktober 2026" sonst hinter dem
%% "bis" und läse sich wie zwei Angaben.
\begin{minipage}[t]{0.44\linewidth}
  \raggedleft\small\color{gedaempft}
  \ifblank{<<GESCHAEFTSFUEHRUNG>>}{}{Geschäftsführung\zeilenwechsel{}<<GESCHAEFTSFUEHRUNG>>\blocktrenner{}}
  \kopfangabe{Kundennr.:}{<<KUNDENNUMMER>>}
  \kopfangabe{Unsere Auftragsnr.:}{<<AUFTRAGSNUMMER>>}
  \kopfangabe{<<ZEITRAUM_BESCHRIFTUNG>>:}{\zusammen{<<LEISTUNGSZEITRAUM>>}}
  \kopfangabe{Abschlag vereinnahmt am}{<<VEREINNAHMT_AM>>}
  \kopfangabe{Ihre Referenz:}{<<KUNDENREFERENZ>>}
  \kopfangabe{Zahlbar ohne Abzug bis}{<<ZAHLUNGSZIEL>>}
\end{minipage}

\vspace{14mm}

%% ---------------------------------------------------------- Titel (Blickfang)
%% Das einzige Flächenelement der Seite. Nummer links, Datum rechts – wie auf
%% der bisherigen Rechnung, damit beide Angaben zusammen gelesen werden.
%% \fboxsep steuert den Innenabstand des \colorbox – waagerecht wie senkrecht.
%% Nachgemessen: Der Balken der bisherigen Rechnung ist 6,0 mm hoch.
%%
%% \raisebox erzwingt Höhe = Ziffernhöhe und Tiefe = 0. Ohne das richtet sich
%% die Boxhöhe nach dem tatsächlichen Text: Die Unterlänge des "g" in "August"
%% macht den Balken unten dicker als oben. Jetzt liegt die Unterlänge im
%% Innenabstand, und die Ziffern haben oben wie unten denselben Abstand.
{\setlength{\fboxsep}{1.4mm}%
\settoheight{\balkenhoehe}{\large\bfseries 0}%
\noindent\colorbox{akzent}{%
  \raisebox{0pt}[\balkenhoehe][0pt]{%
    \makebox[\dimexpr\linewidth-2\fboxsep\relax][s]{%
      \large\bfseries\color{white}<<DOKUMENTART>> <<RECHNUNGSNUMMER>>\hfill <<RECHNUNGSDATUM>>}}}}

\vspace{9mm}

%% --------------------------------------------------------------- Positionen
\begin{tabularx}{\linewidth}{@{}>{\raggedright\arraybackslash}X r@{\hspace{1.5mm}}l r r@{}}
  %% \mbox um jeden Spaltenkopf: In der X-Spalte ist der Inhalt sonst ein Absatz,
  %% dessen erste Zeile tiefer sitzt als die Grundlinie der r-/l-Spalten.
  \mbox{\rubrik{Beschreibung}} & \mbox{\rubrik{Menge}} & & \mbox{\rubrik{Einzelpreis}} & \mbox{\rubrik{Betrag}} \\[1.5mm]
  \midrule
<<POSITIONEN>>
\end{tabularx}

\vspace{5mm}

%% ------------------------------------------------------------------ Summen
\hfill
\begin{minipage}{0.55\linewidth}
  \begin{tabularx}{\linewidth}{@{}>{\raggedright\arraybackslash}X r@{}}
<<SUMMEN>>
  \end{tabularx}
\end{minipage}

%% Der Pflichtsatz der Rechnungsart (siehe \hinweisabsatz oben). Bei jeder
%% anderen Art bleibt der Platzhalter leer und nimmt dann auch keinen Raum ein.
<<ARTHINWEIS>>

%% Begründungen für Positionen ohne Umsatzsteuer (BT-120). Bleibt der Platz-
%% halter leer, entsteht hier nichts – der Abstand darunter trägt beide Fälle.
\vspace{6mm}

<<STEUERHINWEISE>>

%% Die Absetzung der vereinnahmten Abschläge in einer Schlussrechnung – unter
%% dem Gesamtbetrag, nie in der Positionstabelle (ADR-0002), und nach den
%% Steuerhinweisen: Die begründen die Steuerzeilen darüber und gehören zu
%% ihnen. Leer bei jeder anderen Rechnungsart; den Abstand bringt
%% \absetzungskopf mit.
%%
%% Dahinter das Skonto als \hinweisabsatz: nach der Absetzung, weil es das
%% betrifft, was zu zahlen ist. Auf derselben Zeile, damit eine Rechnung ohne
%% Skonto kein Zeichen mehr im .tex hat als vorher.
<<ABSETZUNG>><<SKONTO>>

\vspace{5mm}

%% ------------------------------------------------------------------ Schluss
%% Ohne Zahlungsaufforderung: Das Zahlungsziel steht oben rechts, ein Satz
%% darüber wäre die zweite Nennung desselben Datums. Auf welches Konto zu
%% zahlen ist, sagt die Bankverbindung darunter.
<<SCHLUSSSATZ>>

\vspace{5mm}

Mit freundlichen Grüßen\\[7mm]
<<INHABER>>

%% ---------------------------------------------------------- Bankverbindung
%% Im Textfluss, linksbündig, mit schlichter Beschriftung – wie auf der
%% bisherigen Rechnung. Kein Kasten, keine Linie, keine Versalien. Ob sie
%% dasteht, entscheidet der Renderer (siehe \bankblock oben).
<<BANKVERBINDUNG>>

\vfill

%% ------------------------------------------------------------------ Fußzeile
%% Zwei zentrierte Zeilen ohne Trennlinie darüber, wie auf der bisherigen
%% Rechnung: erst die gesellschaftsrechtlichen Angaben (§35a GmbHG),
%% dann die steuerlichen und die Bankverbindung.
%% Beide Zeilen passen in den Satzspiegel (nachgemessen 152 mm bzw. 138 mm bei
%% 165 mm Breite), deshalb ohne den früheren Trick mit negativem \hspace.
%%
%% Bewusst OHNE Steuernummer: §14 Abs. 4 Nr. 2 UStG verlangt die Steuernummer
%% ODER die USt-IdNr., nicht beides. Mit beiden wurde die zweite Zeile 178 mm
%% breit und passte nicht mehr in die 165 mm Textbreite.
{\centering\scriptsize\color{gedaempft}
<<FIRMA>> · <<STRASSE>>, <<PLZ>> <<ORT>>\ifblank{<<REGISTERGERICHT>>}{}{ · Registergericht <<REGISTERGERICHT>>\ifblank{<<REGISTERNUMMER>>}{}{, <<REGISTERNUMMER>>}}\\
USt-IdNr. <<USTID>> · <<INSTITUT>> · BIC <<BIC>> · IBAN <<IBAN>>\par}

\end{document}
